\section{Detailed numerical verification}
\label{sec:supp-numerics}% Generated by verify_publication_tables.py; do not hand-edit.
\newcommand{\numBoundaryDirection}{1.224744871327}
\newcommand{\numBoundaryDirectionLimit}{1.224744871392}
\newcommand{\numBoundaryLimit}{1.2599210499}
\newcommand{\numBoundaryRatio}{1.2599210420}
\newcommand{\numCenterDistance}{0.040996138377}
\newcommand{\numCompactDirection}{0.6373774341}
\newcommand{\numCompactDirectionLimit}{0.6373774392}
\newcommand{\numCompactMultiplier}{-0.3749999965}
\newcommand{\numCorankEta}{0.715080068980}
\newcommand{\numCorankKrOne}{0.581665969555}
\newcommand{\numCorankKrTwo}{0.336668060890}
\newcommand{\numCorankSlopeFour}{-0.7150801}
\newcommand{\numHuberBoundaryRatio}{1.9999999500}
\newcommand{\numNonDiagonalCosineFive}{0.999999999713}
\newcommand{\numNonDiagonalCosineSix}{0.999999999997}
\newcommand{\numNonDiagonalDirectionFive}{0.671768177074}
\newcommand{\numNonDiagonalDirectionSix}{0.671774660292}
\newcommand{\numNonDiagonalEta}{0.644812526657}
\newcommand{\numNonDiagonalFiniteTOneFive}{0.701047395444}
\newcommand{\numNonDiagonalFiniteTOneSix}{0.701052493549}
\newcommand{\numNonDiagonalMR}{0.6717753807}
\newcommand{\numNonDiagonalSlopeFive}{-0.644805516183}
\newcommand{\numNonDiagonalSlopeSix}{-0.644811825605}
\newcommand{\numNonDiagonalTOne}{0.7010530600}
\newcommand{\numPowerThreeLimit}{1.2778862085}
\newcommand{\numPowerThreeRatio}{1.2778860451}
\newcommand{\numQuarticFour}{-0.124998347562}
\newcommand{\numQuarticOne}{0.004629629625}
\newcommand{\numQuarticThree}{0.374999996875}
\newcommand{\numQuarticTwo}{0.023437499951}
\newcommand{\numRegularLimit}{-1.7777777778}
\newcommand{\numRegularRatio}{-1.7777777766}
\newcommand{\numStrictLimit}{-0.3227486122}
\newcommand{\numStrictRatio}{-0.3227486102}


The local laws of Section~\ref{sec:rates} predict constants, not merely
exponents.  The file \nolinkurl{verify_publication_tables.py} generates every
numerical table and the displayed decimal constants in this section from one
machine-readable snapshot.  It uses 80-digit mpmath arithmetic (100-digit
Decimal arithmetic for the moving family), and repeats the non-diagonal
calculation at 100 digits.  The repeated ratios agree within \(10^{-60}\)
and its intrinsic constants within \(10^{-65}\); we report fewer digits.
These are working precisions and numerical cross-checks, not validated
error bounds.  Bisection and explicit spectral formulas avoid float64
subtraction of nearby multipliers or directions.  The full suite also retains
independent float64 regression tests; no assertion is silently disabled under
\texttt{python -O}.  The following commands run the suite and compare the
stored JSON/LaTeX snapshots, respectively:
\begin{quote}\small\ttfamily
python verify\_all.py\\
python verify\_publication\_tables.py --check
\end{quote}
Omit \texttt{--check} from the second command to regenerate the snapshots.

For the first group, \(\Theta=e_1e_1^\top\) and
\(G=\bigl(\begin{smallmatrix}0&b_1\\b_2&c\end{smallmatrix}\bigr)\).

For the regular row, \(b_1=b_2=1\), \(c=1/2\), and
\(\lambda^\star=1/2\); the ratio is
\((\lambda_\delta-\lambda^\star)/\delta^2\).  For the strict row,
\(b_1=b_2=1/2\), \(c=1\), and \(\lambda^\star=1/4\); the ratio is
\((\lambda_\delta-\lambda^\star)/\delta\).  For the boundary row,
\(b_1=b_2=c=1\), \(\lambda^\star=1\), and the ratio is
\((\lambda^\star-\lambda_\delta)/\delta^{2/3}\).
\begin{center}
\begin{tabular}{@{}llll@{}}
\toprule
Regime & \(\delta\) & Normalized displacement & Limit\\
\midrule
regular & \(10^{-4}\) & \(-1.7777776645\) & \(-1.7777777778\)\\
regular & \(10^{-5}\) & \(-1.7777777766\) & \(-1.7777777778\)\\
strict & \(10^{-6}\) & \(-0.3227484177\) & \(-0.3227486122\)\\
strict & \(10^{-8}\) & \(-0.3227486102\) & \(-0.3227486122\)\\
boundary & \(10^{-9}\) & \(1.2599202562\) & \(1.2599210499\)\\
boundary & \(10^{-12}\) & \(1.2599210420\) & \(1.2599210499\)\\
\bottomrule
\end{tabular}
\end{center}


For the same boundary data, standard Huber gives
\((\lambda^\star-\lambda_\delta^{\rm hub})/\delta=\numHuberBoundaryRatio\) at
\(\delta=10^{-8}\), versus \(1/\beta=2\).  The scaled-Huber regression in
the code also confirms that moving the first saturation threshold to
\(L=1/2\) changes the constant to \(L/\beta=1\), as required by
\eqref{eq:saturated-boundary-rate-main}.

For the admissible profile \(\psi'(x)=x/(1+x^3)^{1/3}\), \(x\ge0\),
we divide the displacement by \(\delta^{3/4}\); for Huber and scaled
Huber we divide it by \(\delta\).  Thus the complete profile table tests
the tail-dependent exponent as well as the leading constant.
\begin{center}
\begin{tabular}{@{}llll@{}}
\toprule
Profile & \(\delta\) & Normalized displacement & Limit\\
\midrule
Huber & \(10^{-4}\) & \(1.999500189915\) & \(2.000000000000\)\\
Huber & \(10^{-6}\) & \(1.999995000019\) & \(2.000000000000\)\\
Huber & \(10^{-8}\) & \(1.999999950000\) & \(2.000000000000\)\\
\(p=3\) & \(10^{-3}\) & \(1.272747989329\) & \(1.277886208493\)\\
\(p=3\) & \(10^{-6}\) & \(1.277857161983\) & \(1.277886208493\)\\
\(p=3\) & \(10^{-9}\) & \(1.277886045147\) & \(1.277886208493\)\\
\(p=3\) & \(10^{-12}\) & \(1.277886207574\) & \(1.277886208493\)\\
\(L=1/2\) & \(10^{-8}\) & \(0.999999987500\) & \(1.000000000000\)\\
\bottomrule
\end{tabular}
\end{center}


The preceding instances are \(2\times2\).  To exercise the general
completion and arbitrary-corank strict law, take
\[
 A^\star=\diag(1,1,0,0),\qquad
 \Theta=\diag\!\left(\tfrac38,\tfrac38,1,\tfrac12\right),\qquad
 \lambda^\star=1,\qquad G=A^\star-\Theta.
\]
Then \(a=3/4\), \(C=\diag(1,1/2)\), and \(\beta=3/2\).  Hard
water-filling gives
\(\eta_F=3/5\),
\(K_F=-\diag(3/5,3/10)\), and
\(\|K_F\|_F=\sqrt{0.45}\).  The smooth-center equation gives
\[
 \eta_R=\numCorankEta,\qquad
 K_R=-\diag(\numCorankKrOne,\numCorankKrTwo),
\]
with \(\|K_F-K_R\|_F=\numCenterDistance\).  Writing
\(\Phi_F^\star\) and \(\Phi_R^\star\) for the two corresponding completed
certificates, the actual smoothing path gives
\begin{center}
\begin{tabular}{@{}llll@{}}
\toprule
\(\delta\) & \(\lambda_\delta-\lambda^\star\) & \(\|\Phi_\delta-\Phi_R^\star\|_F\) & \(\|\Phi_\delta-\Phi_F^\star\|_F\)\\
\midrule
\(10^{-2}\) & \(-7.150296\times10^{-3}\) & \(7.896227\times10^{-5}\) & \(4.098952\times10^{-2}\)\\
\(10^{-3}\) & \(-7.150796\times10^{-4}\) & \(7.858742\times10^{-7}\) & \(4.099607\times10^{-2}\)\\
\(10^{-4}\) & \(-7.150801\times10^{-5}\) & \(7.854955\times10^{-9}\) & \(4.099614\times10^{-2}\)\\
\(10^{-5}\) & \(-7.150801\times10^{-6}\) & \(7.854576\times10^{-11}\) & \(4.099614\times10^{-2}\)\\
\bottomrule
\end{tabular}
\end{center}

Thus the path converges to the smooth rather than the minimum-Frobenius
center, while
\((\lambda_\delta-\lambda^\star)/\delta\to-\eta_R\), as predicted by
Theorems~\ref{thm:general-fenchel-center} and~\ref{thm:rate-strict-general}
and Corollary~\ref{cor:center-agreement-main}.
On the same line, the standard Huber path gives slope
\(-\eta_F=-3/5\) and converges to \(\Phi_F^\star\) to below
\(10^{-11}\) at \(\delta=10^{-5}\), verifying
Proposition~\ref{prop:smoothing-family-main} on the same nonsingleton slice.

The preceding corank-two line is diagonal: the relevant singular dyads are
fixed.  In the notation of Proposition~\ref{prop:strict-constants-closed-main},
both \(P_{\rm big}(s)\) and \(\widetilde B(s)\) are constant.  Hence
\(\partial_s\mathcal W=0\), which proves \(t_1=0\) and
\(\mathcal M_R=0\) rather than merely observing their cancellation.  To test a nonzero first-order
coefficient without random data, take
\[
 A^\star=\diag(1,3/5,0,0),\qquad \lambda^\star=1,
\]
\[
 \Theta=
 \begin{pmatrix}
  1/6&0&2/5&0\\
  0&1/6&0&3/10\\
  1/5&0&2/3&0\\
  0&1/10&0&1/3
 \end{pmatrix},
 \qquad G=A^\star-\Theta.
\]
In the displayed coordinates, \(a=1/3\),
\(C=\diag(2/3,1/3)\), \(\beta=1\), and
\(\eta_R=\numNonDiagonalEta\).  At \(\delta=10^{-5}\), the multiplier ratio is
\[
 (\lambda_\delta-\lambda^\star)/\delta=\numNonDiagonalSlopeFive.
\]
The direction ratio is \(\numNonDiagonalDirectionFive\), already exhibiting
a nonzero linear direction term.  Direct evaluation of
\eqref{eq:t1-closed-main}--\eqref{eq:MR-closed-main} gives
\[
 t_1=\numNonDiagonalTOne,
 \qquad \|\mathcal M_R\|_F=\numNonDiagonalMR.
\]
At \(\delta=10^{-6}\), the observed direction ratio is
\(\numNonDiagonalDirectionSix\), and the cosine between
\((\Phi_\delta-\Phi_R^\star)/\delta\) and \(\mathcal M_R\) equals one to
ten displayed digits (its computed value is \(\numNonDiagonalCosineSix\)).  The finite-scale estimate of
\(t_1\) at \(\delta=10^{-5}\) is \(\numNonDiagonalFiniteTOneFive\).  Thus the table verifies the
predicted vector coefficient, not only its exponent.  Writing
\(t_\delta=(\lambda_\delta-1)/\delta\), the complete sequence is
\begin{center}
\begin{tabular}{@{}llll@{}}
\toprule
\(\delta\) & \(t_\delta\) & \(\|\Phi_\delta-\Phi_R^\star\|_F/\delta\) & \((t_\delta-t_0)/\delta\)\\
\midrule
\(10^{-2}\) & \(-0.637857238093\) & \(0.664828202840\) & \(0.695528856469\)\\
\(10^{-3}\) & \(-0.644112038622\) & \(0.671057559126\) & \(0.700488035110\)\\
\(10^{-4}\) & \(-0.644742427015\) & \(0.671703367737\) & \(0.700996427320\)\\
\(10^{-5}\) & \(-0.644805516183\) & \(0.671768177074\) & \(0.701047395444\)\\
\(10^{-6}\) & \(-0.644811825605\) & \(0.671774660292\) & \(0.701052493549\)\\
\bottomrule
\end{tabular}
\end{center}

This demonstrates
sharpness for one fixed instance; no generic measure statement is inferred.

The compact-root branch of
Corollary~\ref{cor:compact-root-superconvergence-main} is sharp as well.
For \(A^\star=\diag(2,1,0)\), \(\Theta=\diag(1,-1,1)\), and
\(G=A^\star-\Theta\), one has \(a=0\), \(C=(1)\),
\(J_\star=\diag(1/4,1,0)\), and \(\chi_\star=-3/4\).  Thus the predicted
coefficients are \(-3/8\) and
\(\|\mathcal N_0\|_F=\sqrt{26}/8=\numCompactDirectionLimit\ldots\).
At \(\delta=10^{-4}\), high-precision bisection gives the normalized pair
\((\numCompactMultiplier,\,\numCompactDirection)\); the assertion script checks monotone
convergence over four scales.  Thus, even when the compact-polar equation
already has a root, the nearby smooth path has nonzero errors beginning at
cubic and quadratic order.

The doubly cancelled branch in
\eqref{eq:compact-root-quintic-main} uses
\[
 A^\star=\diag(3,2,1,0),\qquad \lambda^\star=1,
\]
\[
 \Theta=\begin{pmatrix}
  27/10&0&0&3/10\\
  0&-16/5&0&1/10\\
  0&0&1/2&1/4\\
  1/5&-3/20&1/10&1
 \end{pmatrix},\qquad G=A^\star-\Theta.
\]
Here \(a=0\), \(C=(1)\), \(d_0=1\),
\(\chi_\star=0\), and \(\chi_2=1/3\).  Hence the predicted quintic
coefficient is \(-1/8\), and the quartic direction coefficient after
subtracting \(-\delta^2J_\star/2\) is
\[
 \diag\!\left(\frac1{216},\frac3{128},\frac38,-\frac18\right).
\]
Extended-precision bisection gives
\begin{center}
\begin{tabular}{@{}ll@{}}
\toprule
\(\delta\) & \((\lambda_\delta-1)/\delta^5\)\\
\midrule
\(10^{-2}\) & \(-0.1248207361\)\\
\(10^{-3}\) & \(-0.1249833498\)\\
\(10^{-4}\) & \(-0.1249983476\)\\
\(10^{-5}\) & \(-0.1249998349\)\\
\bottomrule
\end{tabular}
\end{center}

At \(\delta=10^{-4}\), the diagonal of the normalized quartic direction is
\[
 (\numQuarticOne,\,\numQuarticTwo,\,\numQuarticThree,\,\numQuarticFour),
\]
while the Frobenius norm of its off-diagonal part is below
\(4\times10^{-6}\).  The high-precision assertion script checks all four
scales and the full matrix coefficient.

For Theorem~\ref{thm:crossover-main} we use the one-parameter family
\(b_1=b_2=1\), \(c=1+\varepsilon\), for which
\(\beta_0=\kappa_0=\tfrac12\) and
\(\Delta_\varepsilon=(c^2-1)/(1+c^2)\).  Choosing \(\varepsilon=
\varepsilon(\delta)\) so that \(\Delta_\delta=\zeta\delta^{2/3}\) places the
family in the critical window \eqref{eq:crossover-critical-main}, whose
limit solves \(\tfrac12y^3+\zeta y^2=1\).

\begin{center}
\begin{tabular}{@{}llll@{}}
\toprule
\(\zeta\) & \(\delta\) & \(x_\delta/\delta^{2/3}\) & \(y_\zeta\)\\
\midrule
0 & \(10^{-9}\) & \(1.259920256195\) & \(1.259921049895\)\\
1 & \(10^{-9}\) & \(0.839284832515\) & \(0.839286755214\)\\
3 & \(10^{-9}\) & \(0.552471210127\) & \(0.552474610531\)\\
10 & \(10^{-9}\) & \(0.313769660579\) & \(0.313775961214\)\\
\bottomrule
\end{tabular}
\end{center}


For a valid strict-side moving path, take
\(\Delta_\delta=\delta^{1/3}\); then
\(\Delta_\delta\to0\) while
\(\Delta_\delta/\delta^{2/3}\to\infty\).  The normalization in
\eqref{eq:crossover-inner-main}, whose limit is one for this family, gives
\begin{center}
\begin{tabular}{@{}lll@{}}
\toprule
\(\delta\) & \(x_\delta\sqrt{\Delta_\delta}/\delta\) & \(x_\delta/\widehat x_\delta\)\\
\midrule
\(10^{-6}\) & \(0.979950781607\) & \(0.985083206619\)\\
\(10^{-9}\) & \(0.997994100405\) & \(0.998500862499\)\\
\(10^{-12}\) & \(0.999799770023\) & \(0.999850008712\)\\
\(10^{-15}\) & \(0.999979992294\) & \(0.999985000087\)\\
\bottomrule
\end{tabular}
\end{center}

This path respects the moving-family quantifier.  By contrast, fixing
\(\varepsilon\) and sending \(\delta\) to zero tests the fixed strict-kink
law, not the crossover theorem.

Finally, the oracle returns a matrix direction, not only a multiplier.  At
the boundary instance \(b_1=b_2=c=1\), direct differentiation gives
\[
 \|\mathcal M_0\|_F=\sqrt{3/2},
 \qquad
 \frac{\|\Phi_\delta-\Phi^\star\|_F}{x_\delta}
 \longrightarrow\sqrt{3/2}.
\]
At \(\delta=10^{-15}\), the computed ratio is
\(\numBoundaryDirection\), versus
\(\sqrt{3/2}=\numBoundaryDirectionLimit\).  The same verification includes critical,
strict-side, boundary, and deliberately oscillatory paths for the unified
ratio \(x_\delta/\widehat x_\delta\).

For completeness, the oscillatory sequence is
\[ q_k=10^{-k},\qquad \delta_k=q_k^3,\qquad
\Delta_k=\begin{cases}
0,& k\equiv0\pmod3,\\
3q_k^2,& k\equiv1\pmod3,\\
q_k,& k\equiv2\pmod3.
\end{cases} \]
It alternates among boundary, critical, and strict-side windows, so the
unified ratio has a limit even though the regime changes along the sequence.
The displayed ones are rounded values, not exact identities.
\begin{center}
\begin{tabular}{@{}llll@{}}
\toprule
\(k\) & Window & \(\delta_k\) & \(x_{\delta_k}/\widehat x_{\delta_k}\)\\
\midrule
2 & strict & \(10^{-6}\) & \(0.985083206619\)\\
3 & boundary & \(10^{-9}\) & \(0.999999370040\)\\
4 & critical & \(10^{-12}\) & \(0.999999952844\)\\
5 & strict & \(10^{-15}\) & \(0.999985000087\)\\
6 & boundary & \(10^{-18}\) & \(0.999999999999\)\\
7 & critical & \(10^{-21}\) & \(1.000000000000\)\\
8 & strict & \(10^{-24}\) & \(0.999999985000\)\\
9 & boundary & \(10^{-27}\) & \(1.000000000000\)\\
10 & critical & \(10^{-30}\) & \(1.000000000000\)\\
\bottomrule
\end{tabular}
\end{center}


The separate \texttt{verify\_certificate\_example.py} illustrates why a
certificate must be assessed on the actual returned matrix.  For
\(G=\left(\begin{smallmatrix}0&1\\1&2\end{smallmatrix}\right)\) and
\(\Theta=e_1e_1^\top\), the exact completed direction is
\(\left(\begin{smallmatrix}0&1/2\\1/2&3/4\end{smallmatrix}\right)\).
It is tangent and optimal, whereas the compact polar factor at the same
multiplier has tangency residual \(1/5\).  The script also compares an
eight-step float64 polynomial map using the published Newton--Schulz
coefficients from Xie et al.~\cite{xie2026controlled}.  That diagnostic is not the production BF16 implementation,
a Fenchel-path assertion, or a GPU-training experiment; its floating-point
gap is not an interval-certified bound.

We emphasize what these tables do and do not support.  They confirm the
exponents and leading constants of the local laws on instances where those
constants are known in closed form.  They are not evidence about how often
rank-deficient crossings occur in any application, and no such claim is made.
